\documentclass[10pt]{article}
\usepackage[margin=0.85in]{geometry}
\usepackage{booktabs}
\usepackage{hyperref}
\usepackage{xcolor}
\usepackage{float}
\usepackage{caption}
\usepackage{array}
\usepackage{amsmath}
\usepackage{tikz}
\hypersetup{colorlinks=true,linkcolor=blue!55!black,citecolor=blue!55!black,urlcolor=blue!55!black}
\definecolor{streetone}{HTML}{B08D4F}
\definecolor{streettwo}{HTML}{4C8DFF}
\definecolor{streetthree}{HTML}{16408F}
\definecolor{noise}{HTML}{DDE2EC}
\captionsetup{font=small,labelfont=bf}
\input{generated/macros.tex}
\makeatletter\let\tabinput\@@input\makeatother

\title{How Low to Aim in 2-7 Triple Draw:\\
Pat, Break and Discard Decisions of the Strongest Arena Bots}
\author{Nuttakit Kundum}
\date{October 2026}

\begin{document}
\maketitle

\begin{abstract}
We are building a precomputed solution library for fixed-limit 2-7 triple draw. The betting is
small enough for counterfactual regret minimization (CFR), but each draw adds a second decision:
stand pat or draw, and if drawing, which cards to throw. Before fixing the solver's draw abstraction
we asked how the strongest available programs make that decision. We downloaded \NHands\ hands from
\NMatchesUsed\ matches on the MixedSolver Arena, where every hole card, discard and drawn card is
published, and studied the \NElite\ draw decisions made by the top heads-up and six-max bots.
Four regularities stand out. First, once the number of cards to draw is fixed, the cards to keep
follow one rule (keep the lowest distinct ranks) in at least \RuleLow\% of \RuleN\ draws; the one
real exception is breaking a dealt straight. Second, the drawing target is an eight or better: a one-card draw keeps
a nine as its top card in only \KeptNineHU\% of heads-up and \KeptNineSix\% of six-max draws. Third,
the worst hand they stand pat with falls from an eight on the first draw to a nine or ten on the
last, and moves by two full ranks with the number of cards the opponent just drew: heads-up on the
last draw they always stand pat with a ten against a one-card draw but stand pat with it only
\PatTenVsPat\% of the time against a pat hand. Fourth, the decision to break a made nine, ten or
jack depends on what is left after breaking: T-9 stands pat \PatTNineTwo\% of the time on the second
draw, T-8 only \PatTEightTwo\%, and T-8 with many live outs almost never. We turn these into seven
testable hypotheses for the solver and a concrete draw abstraction: CFR chooses how many cards to
draw, and a rule chooses which.
\end{abstract}

\section{Why the draw needs its own answer}

In fixed-limit 2-7 triple draw each player is dealt five cards, there are four betting rounds with
a cap of five bets, and before each of the last three rounds every live player may exchange any
number of cards. The lowest hand wins: aces are high, straights and flushes count against you, and
the best hand is 7-5-4-3-2 offsuit. The betting tree alone is small. The draws are what make the
game hard to solve: a player holding five cards can keep any of $2^5=32$ subsets, and the right
choice depends on the opponent's range, how many cards the opponent drew, and which cards are
already out of the deck.

Our plan is to let CFR handle the betting and the decision to stand pat or draw, and to choose the
discards by rule. That plan only works if the rule loses little, and it leaves open the question
this paper is mainly about: \emph{how low should a draw aim}? Should a player with a made ten stand
pat or break it to draw to an eight? Is a nine-low draw worth keeping? We do not answer these by
solving here. We look at what strong programs do and write down what a solver should be able to
reproduce.

\section{Data}

The MixedSolver Arena (\url{https://arena.sorawit.dev/competitions}) runs reproducible bot
matches. It holds \NMatchesAll\ completed 2-7 triple draw matches, most of them 10{,}000 to
300{,}000 hands in duplicate mode. For each match the public API exposes about 100 \emph{sample}
hands, which are the first hands of the match and so an unbiased draw, and about 100
\emph{biggest} pots. Every exposed hand shows all hole cards, every action, each player's discards
and the cards drawn.

\paragraph{Choosing the bots.} We ranked every bot by its hand-weighted win rate in matches without
baseline bots (random, always-call and similar). Table~\ref{tab:ranking} lists the top of each
format. We call the heads-up bots from \texttt{paul-gandalf200-4bit} down to
\texttt{2-7-tourmaline-1} and the six-max \texttt{swit-27td-ring*} and \texttt{paul-sauron*}
families \emph{elite}. We kept up to 12 recent matches for each set of elite bots that met, as long
as no baseline bot was seated, and downloaded every exposed hand: \NHands\ hands, of which
\NSampleHands\ are sample hands. Only decisions made by elite bots are counted below
(\NEliteHU\ heads-up, \NEliteSix\ six-max); their opponents in some six-max matches are weaker
bots.

\begin{table}[H]
\centering\small
\caption{Strongest triple draw bots on the arena, win rate in BB/100 as reported by the arena,
weighted by hands, against non-baseline opponents only.}
\label{tab:ranking}
\begin{tabular}{lrrr}
\toprule
Bot & BB/100 & Matches & Hands\\
\midrule
\tabinput generated/ranking.tex 
\bottomrule
\end{tabular}
\end{table}

\paragraph{Parsing.} For each player we replay the hand: start from the dealt cards, and at each
draw record the hand before the draw, the cards kept, the number drawn, how many cards each opponent
who drew earlier on the same round took, and the player's own earlier discards. We wrote our own
2-7 evaluator and checked it against the arena's awarded pots: it picks the same winner in all
\NShowdowns\ showdowns.

\paragraph{Hand classes.} A made hand is five distinct ranks with no straight and no flush. We name
it by its top card (9, T, J) and, for sevens and eights, by its top two cards (7-6, 8-7). So
``a made ten'' is any T-x-x-x-x low.

\section{Which cards to throw}

Given the number of cards drawn, we tested one rule: keep the cards of distinct rank with the lowest
ranks, compared from the top card down, and on a tie keep more suits. Table~\ref{tab:rule} shows
how often the elite bots' kept cards match it.

\begin{table}[H]
\centering\small
\caption{Share of elite draws (\%) whose kept ranks equal the rule ``keep the lowest distinct
ranks'', by number of cards drawn. Counts in brackets.}
\label{tab:rule}
\begin{tabular}{lll}
\toprule
Cards drawn & Heads-up & Six-max\\
\midrule
\tabinput generated/rule.tex 
\bottomrule
\end{tabular}
\end{table}

\StraightMissShare\% of the misses are one decision: breaking a dealt straight such as 7-6-5-4-3
or 8-7-6-5-4 by one card. The rule drops the top card and keeps a four-card run (6-5-4-3), which a
7 or a 2 fills back into a straight. The bots do that in \StraightTop\% of the \StraightN\ cases
and in the other \StraightInside\% drop an inside card instead (keeping 7-5-4-3), which has no
straight risk but a higher top card. So this is a mixed choice between two one-card draws, not a
fixed rule. With any other hand, four-card runs are kept as the rule says. Apart from that one
spot, the question of \emph{which} cards to keep is settled by a rule. The strategic question is
\emph{how many}.

\section{How many cards, and how low to aim}

\paragraph{Draw targets.} When the elite bots draw one card, the highest card they keep is eight or
lower in \KeptEightHU\% of heads-up and \KeptEightSix\% of six-max draws. A nine is the top kept
card in only \KeptNineHU\% (heads-up) and \KeptNineSix\% (six-max) of one-card draws, and when
they draw two, the kept cards are eight or lower in \KeptEightTwoHU\% of heads-up draws. The
bots draw to eights and better; a nine is a hand they may keep, not a draw they chase.

\paragraph{Smoothing a rough eight.} With exactly four distinct cards of eight or lower plus junk on
the first draw, the bots draw one to the four-card low almost always. The exception is 8-7 on top
(for example 8-7-4-2): they break the eight and draw two \BreakEightSevenTwo\% of the time.

\paragraph{When to stand pat.} Figure~\ref{fig:pat} and Table~\ref{tab:pat} show how often a made
hand stands pat. Every made eight stands pat from the second draw on. On the first draw heads-up
bots stand pat with 8-7 \PatEightSevenHU\% of the time but six-max bots only \PatEightSevenSix\%.
A nine stands pat \PatNineDrawTwoHU\% of the time on the second draw heads-up and
\PatNineDrawThreeHU\% on the third; a ten \PatTenDrawThreeHU\% and a jack \PatJackDrawThreeHU\% on
the third. Six-max bots are a step tighter: on the second draw they stand pat with a nine only
\PatNineDrawTwoSix\% of the time.

\begin{figure}[H]
\centering
\input{generated/patfig.tex}
\caption{Share of made hands that stand pat, by the hand's top cards and the draw. Solid lines and
dots: heads-up; dashed: six-max.}
\label{fig:pat}
\end{figure}

\begin{table}[H]
\centering\small
\caption{Stand-pat frequency (\%) of elite bots by made hand. ``--'' means fewer than 20 cases.}
\label{tab:pat}
\setlength{\tabcolsep}{4.5pt}
\begin{tabular}{llrrrrrrrrrrr}
\toprule
& & 7-5 & 7-6 & 8-5 & 8-6 & 8-7 & 9 & T & J & Q & K & A\\
\midrule
\tabinput generated/pat.tex 
\bottomrule
\end{tabular}
\end{table}

\paragraph{The opponent's draw moves the line.} Heads-up, the player in position draws second and
sees how many cards the opponent took. Table~\ref{tab:vsdraw} splits the stand-pat frequency by
that number. On the last draw, a ten stands pat \PatTenVsOne\% of the time when the opponent drew
one card, a jack \PatJackVsOne\% and a queen \PatQueenVsOne\%. When the opponent stood pat, the same
ten stands pat only \PatTenVsPat\% of the time and even a nine only \PatNineVsPat\%: against a pat
hand the bots break and draw to beat it. There is no fixed answer to how low to aim; the target is
set by the hand the opponent is likely to show down.

\begin{table}[H]
\centering\small
\caption{Heads-up, second player to draw: stand-pat frequency (\%) by made hand and by how many
cards the opponent just drew. ``--'' means fewer than 10 cases.}
\label{tab:vsdraw}
\begin{tabular}{llrrrrr}
\toprule
Draw & Opponent & 9 & T & J & Q & K\\
\midrule
\tabinput generated/vsdraw.tex 
\bottomrule
\end{tabular}
\end{table}

\paragraph{What is left after breaking matters as much as the hand.} The stand-pat frequency is not
monotone in the second card (Table~\ref{tab:second}). Heads-up on the second draw, T-9 stands pat
\PatTNineTwo\% of the time while T-8 stands pat \PatTEightTwo\%; on the third draw J-9 stands pat
\PatJNineThree\% and J-8 \PatJEightThree\%. Breaking T-8-x-x-x leaves an eight draw; breaking
T-9-x-x-x leaves only a nine draw, which (as above) the bots do not want to chase. A made hand is
compared with the draw it would become, not with a fixed threshold.

\begin{table}[H]
\centering\small
\caption{Stand-pat frequency (\%) by the top two cards of a made nine, ten or jack.}
\label{tab:second}
\setlength{\tabcolsep}{4pt}
\begin{tabular}{llrrrrrrrrrrr}
\toprule
& & 9-6 & 9-7 & 9-8 & T-6 & T-7 & T-8 & T-9 & J-7 & J-8 & J-9 & J-T\\
\midrule
\tabinput generated/second.tex 
\bottomrule
\end{tabular}
\end{table}

\section{Card removal and blockers}

\paragraph{Live outs.} We counted, for each made hand, the \emph{live outs after breaking}: unseen
cards that turn the four lowest cards into a no-pair eight or better, removing cards in the
player's own hand and own earlier discards, and removing ranks that would fill a straight. Within
the same hand class, the bots break far more often when those outs are many
(Table~\ref{tab:outs}). On the second draw, T-8 with nine or more live outs stands pat
\OutsTEightManyTwo\% of the time; with seven or eight it stands pat \OutsTEightFewTwo\%. This is card
removal from the player's own hand, and the bots clearly use it.

\begin{table}[H]
\centering\small
\caption{Heads-up stand-pat frequency (\%) for made nines and tens by live outs after breaking.
Counts in brackets; ``--'' means fewer than 10 cases.}
\label{tab:outs}
\begin{tabular}{llrrr}
\toprule
Draw & Hand & $\le 6$ outs & 7--8 outs & 9+ outs\\
\midrule
\tabinput generated/outs.tex 
\bottomrule
\end{tabular}
\end{table}

\paragraph{Own discards.} A separate question is whether cards a player threw earlier, and so knows
are not in the deck, change the decision. We could not test it: players rarely throw cards that
would later be their outs, so only \DeadSomeN\ made nines and tens on the second or third draw had
any. \paragraph{Opponent blockers.} Holding a low card also lowers the chance that the opponent
makes a strong hand. Hand histories cannot separate this from hand strength, because the cards
that block the opponent are the same cards that make one's own hand good. A solver that tracks
exact cards, rather than grouping hands into buckets, gets both effects without special handling.

\section{What wins at showdown}

Table~\ref{tab:showdown} uses only sample hands and only showdowns between two players. Heads-up,
\SdCumEightSevenHU\% of pots are won by 8-7 or better and \SdCumNineHU\% by a nine or better; in
six-max, \SdCumEightSevenSix\% are won by 8-7 or better. Holding a nine at a heads-up showdown wins
\SdWinNineHU\% of the time, a ten \SdWinTenHU\% and a jack \SdWinJackHU\%; in six-max a jack wins only
\SdWinJackSix\%. These are conditional on reaching showdown, so weak hands that get there are
partly bluff-catchers and partly snows called down, but they agree with the stand-pat lines: a
heads-up nine or ten is close to an even showdown hand, an eight is a strong one, and six-max pots
need a step lower. Elite bots also snow (stand pat with a pair to bluff) on the last draw
\SnowThreeHU\% of the time heads-up and \SnowThreeSix\% six-max when they hold a pair.

\begin{table}[H]
\centering\small
\caption{Two-player showdowns in sample hands (\SdNHU\ heads-up, \SdNSix\ six-max): cumulative share
of pots won by this hand or better, and the win rate of a player holding it.}
\label{tab:showdown}
\setlength{\tabcolsep}{4.5pt}
\begin{tabular}{llrrrrrrrrr}
\toprule
& & 7-6 & 8-6 & 8-7 & 9 & T & J & Q & K & A\\
\midrule
\tabinput generated/showdown.tex 
\bottomrule
\end{tabular}
\end{table}

\section{Do the strong bots agree?}

Table~\ref{tab:bots} gives the hardest cells bot by bot. Heads-up, the elite bots differ by about
twenty points on the nine at the second draw (46--66\%) and by fifteen on the ten at the last draw
(41--55\%), and agree more closely on the jack (21--30\%). Six-max is less settled: on the first draw the two \texttt{swit} bots always stand
pat with a made eight while the \texttt{sauron} bots do so about half the time. That is the cell in
which a solver's answer would be most informative.

\begin{table}[H]
\centering\small
\caption{Stand-pat frequency (\%) of each elite bot with at least 15 cases per cell. D2~9 means a
made nine at the second draw.}
\label{tab:bots}
\begin{tabular}{lrrr}
\toprule
\tabinput generated/bots.tex 
\bottomrule
\end{tabular}
\end{table}

\section{Hypotheses for the solver}

\begin{enumerate}
\item[H1] \textbf{Aim at eight or better.} Drawing hands should keep eight-or-better cards; a
one-card draw to a nine is rarely right heads-up and almost never six-max.
\item[H2] \textbf{The pat line falls by draw.} Stand pat with an eight from the first draw; a nine
becomes a mixed decision on the second draw heads-up and mostly a pat on the third; a ten and a
jack are mixed on the third.
\item[H3] \textbf{The opponent's draw count sets the target.} On the last draw, against a one-card
draw a ten is a pure pat and a jack is mixed; against a pat hand even a ten breaks.
\item[H4] \textbf{Compare with the draw left after breaking.} Break when the remaining four cards
are an eight-or-better draw with many live outs; keep T-9 and J-9 rather than break into a nine
draw.
\item[H5] \textbf{Heads-up showdowns are won around a nine.} A nine or ten is close to an even
showdown hand heads-up and a jack is weak in six-max pots.
\item[H6] \textbf{Snows exist but are rare,} a few percent of paired hands on the last draw and
almost none earlier.
\item[H7] \textbf{Six-max plays a step tighter than heads-up,} with the first-draw 8-7 pat as the
open question.
\end{enumerate}

\section{What this means for the solver design}

The draw decision splits cleanly in two. \emph{Which} cards to keep, given how many, follows a rule
the strong bots break less than one time in a hundred, so the solver can fix it and save the
$2^5$ discard branches. The one place to keep both options open is breaking a dealt straight,
where the top card and an inside card are both played. \emph{How many} cards to draw is strategic, and not only pat against draw:
breaking a rough 8-7 to draw two on the first draw, or breaking a made ten instead of standing pat,
are real choices with mixed frequencies. We therefore let CFR choose the number of cards (in
practice: stand pat, draw the natural number, or break one more card), with the kept cards set by
the rule. Because H4 depends on exact live outs, the solver should carry exact cards, including
each player's own discards, rather than bucket hands by strength.

\section{Limits}

These bots are strong, not proven equilibria; agreement between them is evidence, not proof. Only
about 200 hands per match are published, so cells far from the main line have small counts. The
``biggest'' hands over-represent large pots, which is why the showdown section uses sample hands
only. Showdown win rates are conditional on reaching showdown. Some six-max decisions were made
against weaker opponents than the elite set. Hypotheses H1--H7 are what we will test against our own
solver once it is built.

\paragraph{Reproducing.} \texttt{research/td27\_arena/collect.py} downloads the hands,
\texttt{analyze.py} produces every number, and \texttt{build\_paper.py} writes the tables. Data were
downloaded on 4 October 2026.

\end{document}
